\begin{afterword}
\summaryhead

The post returns to the pulp covers on which a bug-eyed monster carries off a human woman. An artist, Fred, reasons that the monster will prefer human women because it can see that their skin is pleasanter than scales. The author's diagnosis is that Fred treats sexiness, a function of two arguments (an admirer and an entity), as a function of the entity alone.

An equally valid view, the post says, is that each speaker's ``sexy'' names a different one-place function, which need not mention the speaker: Fred uses Sexiness\_20934, the monster Sexiness\_72546, and everyone can agree on the values of both. Currying, a technique from programming, unifies the views: the two-place function, given an admirer, returns a fixed one-place function. The same analysis is applied to Putnam's Twin Earth, where ``water'' can be read as a function that takes a world and returns its watery stuff, or as the fixed result for our world, H20. The post says the debate's confusion came from currying or uncurrying by instinct, and that the analysis will later help to taboo ``objective'', ``subjective'' and ``arbitrary''.

\responsehead

The post's central distinction is correct and clearly shown. Whether a woman is sexy depends on who is looking, and there are two consistent ways to say so: make the admirer an argument of the function, or let each admirer's word name a different fixed function. Currying shows exactly how the two forms relate, and the Fred example makes the error concrete.

The distinction is also well established in philosophy, which supports the post. The two-place form is how relativists state their view, as a ``hidden parameter'' in predicates such as ``right'' and ``beautiful'', and response-dependent theories of value, which define a property by how subjects respond to it, tie properties to subjects in the same way. The one-place form, with a different function for each speaker, is what the literature on moral language calls contextualism. The currying step itself has the structure of David Kaplan's distinction between character and content: a function that takes a context and returns a content. For ``water'', the two readings correspond to the two ways of evaluating a world in two-dimensional semantics, as the notes on ``Heat vs.\ Motion'' explain.

The post's one large claim goes beyond its analysis. It says that ``the whole confusingness of the subsequent philosophical debate'' about Twin Earth came from an instinct to curry or uncurry. No argument is given. The two-level answer had already been proposed, and it has drawn criticism of its own: about a priori knowledge, and about whether the function from worlds to what ``water'' picks out is part of meaning at all.

The post says the analysis is meant for later use on words like ``objective'' and ``subjective'', and the book goes on to apply it to morality. For moral words each form has a known difficulty, lost disagreement for the first and changing tastes for the second; the notes record them, and the later posts are where they apply.

In short: a clear and correct account of two-place words, matching analyses philosophers already use; its one larger claim, that the Twin Earth debate came from an instinct to curry or uncurry, is stated without argument.
\end{afterword}
